In this section we formally define the language and rules of $\LambdaTau$ which are not so different from simply-typed \lambdacalc and the standard reductions rules of pure \lambdacalc.

\subsection{Language}

Like in pure \lambdacalc, expressions in $\LambdaTau$ are called \lambdaterms. The following inductive definition establishes how the set \(\LambdaTau\) of all \lambdaterms is constructed. To start with, we assume the existence of an infinite set \(\mathbb{V}\) of so-called variables:
\[
    \mathbb{V} = \{ x, y, z, \dots \}
\]

Below we define the inductive rules of the language. The language is the language of pure \lambdacalc augmented with a rule to handle typed variables as in \simplyTypedLambdaCalc and the \textit{any} atom written $*$.

\begin{reusable}{definition}{lambdaTypesLanguageRules}
    The set \(\LambdaTau\) of all terms is formally defined as:\blockspacer
    \begin{tabular}{rl}
        \textit{(any)}                  & $*\in \LambdaTau$ \\
        \textit{(variable)}             & If \(u \in \mathbb{V}\) then \(u \in \LambdaTau\). \\
        \textit{(application)}          & If \(M, N \in \LambdaTau\) then \((MN) \in \LambdaTau\).\\
        \textit{(typed abstraction)}    & If \(u \in \mathbb{V}\) and \(T, M \in \LambdaTau\) then \((\tabs u:T.M) \in \LambdaTau\).\\
        \textit{(untyped abstraction)}  & If \(X \in \mathbb{V}\) and \(M \in \LambdaTau\) then \((\uabs X.M) \in \LambdaTau\).\\
    \end{tabular}
\end{reusable}

\begin{reusable}{definition}{definitionType}
    In a typed abstraction $\tabs x:X.M$, the term $X$ is called the type of the variable $x$.
\end{reusable}

\begin{reusable}{definition}{definitionTypedAndUntyped}
    An abstraction with no type, like $\uabs X.M$ is called an \textit{untyped abstraction}.
    An abstraction with a type, like $\tabs x:X.M$, is called a \textit{typed abstraction}.
\end{reusable}

We use the term type to not confuse the reader but we can equivalently see types as set of terms with some constraints defined in the type.

There are two major differences with \simplyTypedLambdaCalc: first, we include a very specific term called \text{any} and written $*$, and second, the type of the variable in an abstraction is itself a \lambdaterm.
We include parenthesis to be able to disambiguate associativity. Therefore, the grammar looks like this:
\[
\mathbb{E} := * \mid \mathbb{V} \mid \mathbb{E}\ \mathbb{E} \mid \lambda \mathbb{V}.\mathbb{E} \mid \lambda \mathbb{V}\oftype\mathbb{E}.\mathbb{E} \mid (\mathbb{E})
\]

We now give some examples of expressions from this language.
{
    \newcommand{\sep}{,\ }

    \begin{reusable}{examples}{examplesValidExpressionInLambdaTau}
    \[
    \begin{array}{ll}
        * \sep a \sep b \sep c \sep X \sep Y \sep Z \sep myvar & \text{\textit{any} and single variables} \\
        \app A B \sep \app A \app B C \sep \app A {(\app B C)} \sep \app {(\app A B)} C & \text{standard applications} \\
        \uabs X.Y \sep \uabs X.X \sep \lambda X.* \sep \lambda X.a & \text{untyped abstractions} \\
        \tabs x:Y.x \sep \tabs x:Y.* \sep \tabs x:*.*, \tabs X:Y.X, \tabs x:Y.a & \text{typed abstractions} \\
        \tabs x:X.\tabs y:Y.\app x y \sep \uabs X.\uabs Y.X \sep \uabs X.\uabs Y.* & \text{abstraction with multiple variables} \\
        \tabs x:(\tabs y:z.y).x \sep \tabs x:\app {(\tabs y:z.y)} z.x & \text{complex types} \\
    \end{array}
    \]
    \end{reusable}
}

In the remainder of this paper, we will generally use lowercase variables for the variables of typed abstractions and uppercase variables for the variables of untyped abstractions. For examples: $\tabs x:X.M$ and $\uabs X.Y$.

\begin{lemma}
    The language of pure \lambdacalc $\Lambda$ is a subset of $\LambdaTau$. More formally, for all $M \in \Lambda$, then $M \in \LambdaTau$, or even shorter, $\Lambda \subseteq \LambdaTau$.
\end{lemma}

\begin{proof}
    The rules for variables, application and untyped abstraction are the rules defining $\Lambda$. For $\LambdaTau$ We extend them with types and \textit{any}. Consequently, any term generated in $\Lambda$ can also be generated by using those three rules.
\end{proof}

\subsection{Bridge with Programming}

At this stage, we can already make some parallels with computer science. Let us review the following concepts that certainly feel familiar to computer scientists. Most concepts that we borrow here are related to functions and will hopefully help build some kind of intuition on how this new language attempts to close the gap between terms and types. 
\begin{itemize}
    \item In a typed lambda abstraction $\tabs x:T.B$, we call $x$ the variable, $T$ the type of the variable $x$ and $B$ is called the body of the abstraction.
    \item In an untyped lambda abstraction $\uabs X.B$, $B$ is the body and the variable $X$ is untyped. It is somehow equivalent to the \texttt{any} type in Typescript, \texttt{Object} in Java and \texttt{any} or \texttt{interface\{\}} in Go. Note that we will add a rule to be able to interchangeably write $\uabs X.B$ and $\uabs X:*.B$. 
    \item The application $\app F x$ can be seen as the application of the function $F$ to $x$. However, it feels natural to developers in typed languages that if $F$ expects a variable of type $T$ as input, then $x$ must be of type $T$ for the function to be \textit{executed}. Otherwise, either of these situations usually happen: the compiler stops with an error, the program crashes when attempting to execute the function, the program raises an execution error, the function has an undefined behavior like returning an arbitrary value.
    \item Functions in typed programming languages and in type theory applied to proof systems also share two common concepts: function signatures and return types. These concepts are so important in order to understand the motivation of this paper that they will be expanded later in their own section.
\end{itemize}

To anchor what we have seen so far, let us review a list of commented examples of expressions in $\LambdaTau$.
Note that in this paper, I attempted to take lower case letters for variables of typed lambda abstractions and upper case letters for the rest.

This language naturally extends the concept of \alphaequivalence defined in pure \lambdacalc to the typed abstraction rule. However, \betareduction is extended to a new reduction operator called \taureduction that we shall introduce very soon. The letter $\tau$ has been chosen arbitrarily because it is often used to represent a type in various papers and \taureduction is really about typed reduction.
